% =====================================================================
%  RESUMEN / ABSTRACT
% =====================================================================

% Rótulo común para los dos bloques (si lo prefieres, muévelo a preambulo.tex)
\newcommand{\rotuloresumen}[1]{%
  \begin{center}
    {\LARGE\bfseries #1\par}
  \end{center}
  \vspace{0.8\baselineskip}%
}

\thispagestyle{plain}
\markboth{}{}
\vspace*{2\baselineskip}

\addcontentsline{toc}{chapter}{Resumen}
\rotuloresumen{Resumen}

\noindent Este trabajo presenta la formalización en el asistente de demostración Lean 4 de la lógica modal $\K$ y de sus dos resultados fundamentales: los teoremas de corrección y completitud respecto a la semántica de Kripke. Frente al estilo conciso de una biblioteca como Mathlib, esta formalización se ha desarrollado con una intención deliberadamente didáctica para hacer visible cada paso del razonamiento lógico. La implementación abarca desde la sintaxis básica hasta la construcción del modelo canónico y una versión del lema de Lindenbaum basada en el lema de Zorn. La memoria expone los fundamentos lógicos y computacionales necesarios sin presuponer conocimientos previos, justificando las decisiones de diseño y las tácticas empleadas, y culmina con un ejemplo ilustrativo de aplicación de la lógica modal.

\vspace{3\baselineskip}

\addcontentsline{toc}{chapter}{Abstract}
\rotuloresumen{Abstract}

\begin{otherlanguage}{english}
\noindent This work presents the formalization, in the Lean 4 proof assistant, of the modal logic $\K$ and of its two fundamental results: the soundness and completeness theorems with respect to Kripke semantics. In contrast with the concise style of a library such as Mathlib, this formalization has been developed with a deliberately didactic aim, so that every step of the logical reasoning remains visible. The implementation ranges from the basic syntax to the construction of the canonical model and a version of Lindenbaum's lemma based on Zorn's lemma. The text sets out the logical and computational background required, presupposing no previous knowledge, justifying the design decisions and the tactics employed, and closes with an illustrative example of an application of modal logic.
\end{otherlanguage}